\chapter{Préface}

Ce livre sert à s'entraîner. Il rassemble des questions d'entretien en finance quantitative et leurs réponses
complètes, pour les chercheurs et traders quantitatifs, les traders, les analystes, les data scientists, les
développeurs quantitatifs, les ingénieurs logiciel et basse latence et les quants de risque, chez les teneurs de
marché, les firmes crypto, les banques et les hedge funds.

Les questions ont été écrites pour ce livre. Elles suivent les thèmes, le niveau et la forme que rapportent
publiquement les candidats, et ce ne sont pas des fuites~: aucune question n'est présentée comme venant d'une firme ou
d'un entretien particuliers.

C'est le premier des trois livres d'entretien, et la plupart des candidats en ont besoin. Il entraîne la vitesse,
les probabilités et les mathématiques. Le livre~19 couvre les marchés, les jeux de trading et le dernier tour, et le
livre~20 le code.

\section*{Comment utiliser ce livre}

Il faut répondre à chaque question à voix haute ou sur papier avant de regarder sa réponse, et toujours finir sur une
conclusion nette. Les réponses sont à la fin du livre, pour qu'on ne les voie pas par accident. Chaque question renvoie
à sa réponse, et chaque réponse renvoie à sa question.

Les étoiles donnent le temps qu'il faut à un candidat préparé~: $\star$ deux ou trois minutes, $\star\star$ environ
sept, $\star\star\star$ environ dix, et $\star\star\star\star$ quinze ou plus. Les étiquettes à côté des étoiles
nomment les métiers et les types de firme où la question a sa place, et disent si elle convient à un candidat junior
ou confirmé. On peut sauter ce qui ne correspond pas à son profil.

Une réponse donne le résultat et la démonstration. Quand cela en vaut la peine, elle donne aussi l'erreur classique,
les relances qu'un recruteur posera probablement, et pourquoi la question compte sur un desk.

Classer les questions par thème trahit la méthode, et un vrai test ne le fait pas. Le livre se termine donc par des
épreuves chronométrées~: douze questions mêlées en deux sections, avec des réponses numériques et sans indice. Il faut
les passer dans des conditions réelles, d'une seule traite. Les exercices de vitesse fonctionnent de même~: lancer un
chronomètre et se noter honnêtement.

\IfFileExists{parts/interviews-1/fr/tracks.tex}{% Generated by tools/make_interview_tracks.py. Do not edit.
\section*{Quels chapitres travailler}
Chaque case compte les questions d'un chapitre étiquetées pour ce métier ou ce type de firme. QR désigne le chercheur quantitatif (quant researcher), QT le trader quantitatif, DS le data scientist, QD le développeur quantitatif, Dev le développeur logiciel, LLE l'ingénieur basse latence, Risque le quant de risque~; TM le teneur de marché, Cr la firme crypto, Bq la banque, HF le hedge fund. Les épreuves et les exercices de vitesse sans étiquette sont pour tout le monde.
\par\medskip\noindent{\footnotesize\setlength{\tabcolsep}{2.5pt}%
\begin{tabular}{@{}>{\raggedright\arraybackslash}p{6.2cm}rrrrrrrrrrrrr@{}}
\hline
Chapitre & QR & QT & Trader & Analyste & DS & QD & Dev & LLE & Risque & TM & Cr & Bq & HF \\
\hline
\hyperref[ch:iv1:mental-technique]{Technique mentale} & 19 & 40 & 34 & 12 & 1 & 5 & 4 & 1 & 15 & 37 & 28 & 33 & 35 \\
\hyperref[ch:iv1:sequences-patterns-and-aptitude]{Suites, régularités et aptitudes} & 4 & 30 & 29 & 25 & 3 & 7 & 6 &  & 6 & 29 & 9 & 27 & 13 \\
\hyperref[ch:iv1:estimation]{Estimation} & 8 & 21 & 17 & 13 & 4 & 9 & 6 & 6 & 8 & 22 & 13 & 16 & 24 \\
\hyperref[ch:iv1:number-riddles]{Énigmes sur les nombres} & 24 & 25 & 17 & 1 &  & 11 & 5 &  &  & 25 & 14 & 14 & 25 \\
\hyperref[ch:iv1:counting-and-symmetry]{Comptage et symétrie} & 40 & 40 & 29 & 7 & 2 & 13 & 1 &  & 6 & 39 & 28 & 31 & 40 \\
\hyperref[ch:iv1:conditional-probability-and-bayes]{Probabilité conditionnelle et Bayes} & 38 & 40 & 24 & 14 & 16 &  &  &  & 18 & 32 & 17 & 24 & 39 \\
\hyperref[ch:iv1:expectation-and-indicators]{Espérance et variables indicatrices} & 25 & 23 & 16 & 3 & 6 & 4 & 4 & 1 & 7 & 25 & 20 & 17 & 25 \\
\hyperref[ch:iv1:waiting-times-and-markov-chains]{Temps d'attente et chaînes de Markov} & 25 & 25 & 15 & 3 & 2 & 1 &  &  & 6 & 22 & 13 & 14 & 25 \\
\hyperref[ch:iv1:random-walks-and-martingales]{Marches aléatoires et martingales} & 22 & 25 & 15 & 4 & 1 & 2 &  &  & 14 & 24 & 12 & 13 & 25 \\
\hyperref[ch:iv1:continuous-and-geometric-probability]{Probabilités continues et géométriques} & 25 & 25 & 16 & 3 & 9 & 5 & 1 & 1 & 3 & 24 & 17 & 18 & 25 \\
\hyperref[ch:iv1:order-statistics]{Statistiques d'ordre} & 15 & 15 & 3 & 4 & 10 & 2 &  & 1 & 4 & 14 & 11 & 10 & 14 \\
\hyperref[ch:iv1:optimal-stopping]{Arrêt optimal} & 20 & 20 & 17 & 5 & 5 &  &  &  &  & 15 & 11 & 11 & 19 \\
\hyperref[ch:iv1:games]{Jeux} & 20 & 20 & 14 & 6 &  & 4 &  &  &  & 20 & 14 & 13 & 19 \\
\hyperref[ch:iv1:logic-riddles]{Énigmes de logique} & 35 & 35 & 18 & 19 & 1 & 14 & 3 &  &  & 33 & 30 & 30 & 35 \\
\hyperref[ch:iv1:statistics]{Statistiques} & 40 & 40 & 5 & 16 & 35 & 1 &  & 1 & 10 & 28 & 24 & 25 & 40 \\
\hyperref[ch:iv1:regression]{Régression} & 30 & 22 & 2 & 4 & 25 & 2 &  &  & 16 & 24 & 20 & 25 & 30 \\
\hyperref[ch:iv1:linear-algebra]{Algèbre linéaire} & 25 & 22 &  &  & 16 & 15 & 1 &  & 10 & 14 & 14 & 16 & 25 \\
\hyperref[ch:iv1:calculus-and-optimisation]{Calcul et optimisation} & 20 & 16 & 3 & 4 & 8 & 9 & 4 & 1 & 1 & 15 & 14 & 14 & 20 \\
\hyperref[ch:iv1:machine-learning]{Apprentissage automatique} & 45 & 12 &  & 5 & 45 & 10 &  &  & 2 & 32 & 31 & 31 & 45 \\
\hyperref[ch:iv1:stochastic-calculus]{Calcul stochastique} & 30 & 30 &  &  &  & 24 &  &  & 18 & 8 & 8 & 23 & 30 \\
\hyperref[ch:iv1:time-series]{Séries temporelles} & 20 & 20 &  &  & 15 &  &  &  & 7 & 9 & 6 & 12 & 20 \\
\hline
\end{tabular}}
}{}

\section*{Contribuer}

La source de ce livre se trouve dans un dépôt git public~:
\begin{center}
\url{https://github.com/bvirrion/one-quant-book}
\end{center}
Les corrections et les meilleures questions sont les bienvenues~; voir \texttt{CONTRIBUTING.md} à la racine du dépôt.
